\begin{table}[t]
\centering
\begin{threeparttable}
\caption{Crawler-access restriction and abnormal filing-window illiquidity}
\label{tab:main}
\small
\begin{tabular}{lcc}
\toprule
 & \multicolumn{2}{c}{Borusyak et al.\ (2024)} \\
\cmidrule(lr){2-3}
 & (1) & (2) \\
\midrule
Crawler block $\times$ Post & $0.0134^{***}$ & $0.0140^{***}$ \\
                            & (0.0052)       & (0.0051) \\
Pre-trend $p$ (12 period)   & 0.48           & 0.52 \\
\midrule
Controls         & No   & Yes \\
Firm FE          & Yes  & Yes \\
Filing-month FE  & Yes  & Yes \\
Observations     & 25{,}516 & 25{,}516 \\
\bottomrule
\end{tabular}
\begin{tablenotes}\footnotesize
\item \textit{Notes.} The dependent variable is the abnormal filing-window Amihud illiquidity (Eq.~\ref{eq:twfe}): the log difference between the $[-1,+1]$ trading-day event window and the $[-60,-10]$ local benchmark. We report the \citet{BorusyakJaravelSpiess2024} imputation estimator as our primary specification, without (column~1) and with (column~2) the control set, each with a 12-period pre-trend test; 580 treated firms and 25{,}516 firm-filings. The control set (log assets, leverage, ROA, market-to-book, tangibility, and analyst coverage, winsorized and within-sample demeaned) barely moves the estimate, because the dependent variable is already an abnormal, within-event difference that nets out slow-moving firm characteristics. The transparent two-way fixed-effects benchmark is very close ($+0.0109$ to $+0.0111$, $p<0.05$) and is reported in the robustness section (Table~\ref{tab:twfe_robustness}); the level form and a window-by-window decomposition are also there (Table~\ref{tab:eventspec}). Standard errors clustered by firm in parentheses. $^{*}$, $^{**}$, $^{***}$ denote significance at the 10\%, 5\%, and 1\% levels.
\end{tablenotes}
\end{threeparttable}
\end{table}
